Recovering temporally consistent 3D geometry from two to four cameras is a core
challenge in computer vision, with direct relevance to robotics and autonomous
manipulation. Spatial foundation models achieve high geometric fidelity by
reconstructing each frame independently, but produce high-frequency jitter when
assembled into a sequence. Geometry-refinement approaches and dedicated 4D
architectures reduce jitter at the cost of geometric accuracy, leaving no prior
method that achieves both simultaneously.

This thesis proposes a decoupled paradigm: apply a strong per-frame spatial
estimator, lift its outputs into a common coordinate frame via a lightweight
temporal alignment strategy, and suppress residual jitter through
Confidence-Weighted Temporal Smoothing (Opt), a training-free Gaussian filter
restricted to static scene regions. Geometric accuracy and temporal consistency
conflict only when temporal constraints penalise dynamic regions where motion is
genuine; decoupling the two stages avoids this trade-off by design.

We benchmark six reconstruction models on
DexYCB and Hi4D under two-, three-, and four-view configurations, using a unified
framework covering Chamfer distance, per-region accuracy and completeness, jitter,
drift, and camera pose error. Opt eliminates high-frequency jitter entirely and
reduces mean jitter by 60--91\,\% across four spatial backbones while leaving
Chamfer distance flat or improved for three of the four. The resulting pipeline
Pareto-dominates native 4D architectures on the geometry--jitter trade-off at
every view count. The remaining open problem is drift, which the low-pass filter
cannot address by design; anchor-constrained optimisation is the natural next step.